\chapter{节奏与节拍}
\label{sec:rhythm-and-meter}

\subsection*{学完这一章能做什么}
\begin{itemize}
  \item 听到一段音乐能数出“一拍一拍”的脉冲，并说出它是 4/4、3/4 还是 6/8。
  \item 看懂音符时值、附点、连音与休止，能把“多长”换算成拍数和秒数。
  \item 能听出并写出基本的切分与反拍，不再把所有音符都摆在正拍上。
  \item 在夜半里把工程速度、拍号、标尺与量化设成同一套口径，用节拍器对齐后再写节奏。
\end{itemize}

\subsection{脉冲、拍与速度}

音乐先有脉冲（pulse），才有“拍”。拍（beat）是被组织起来的等间隔脉冲；把拍分组，就得到小节（bar）。拍号告诉你“每小节几拍”，速度（tempo）告诉你“一拍多快”。

速度用 BPM（beats per minute，每分钟拍数）表示。若一拍就是拍号分母所指的时值，则：

\begin{equation}
  T_{\text{beat}} = \frac{60}{\mathrm{BPM}}\ \text{秒},
  \qquad
  T_{\text{bar}} = N_{\text{beats}} \times T_{\text{beat}}.
  \label{eq:beat-duration}
\end{equation}

例如 120 BPM 时一拍 0.5 秒，4/4 的一小节为 2 秒，3/4 的一小节为 1.5 秒。速度改变的是“多快”，不改变音高与音色；同一段旋律加快一倍，听感更急，但音名不变。

\begin{center}
\begin{tikzpicture}[x=1cm,y=1cm,font=\scriptsize]
  % 4/4：两小节
  \node[anchor=west] at (-0.1,3.6) {\textbf{4/4}（每小节四拍，强–弱–次强–弱）};
  \foreach \i in {0,...,7} {
    \draw[black!25] (\i,2) rectangle ++(1,0.65);
  }
  \fill[blue!28] (0,2) rectangle (1,2.65);
  \fill[blue!16] (2,2) rectangle (3,2.65);
  \fill[blue!28] (4,2) rectangle (5,2.65);
  \fill[blue!16] (6,2) rectangle (7,2.65);
  \node at (0.5,2.33) {强};
  \node at (2.5,2.33) {次强};
  \draw[thick] (4,1.9) -- (4,2.75) node[above] {\tiny 小节线};
  \node[anchor=west] at (8.1,2.33) {两小节};

  % 3/4：两小节
  \node[anchor=west] at (-0.1,1.2) {\textbf{3/4}（每小节三拍，强–弱–弱）};
  \foreach \i in {0,...,5} {
    \draw[black!25] (\i,0) rectangle ++(1,0.65);
  }
  \fill[blue!28] (0,0) rectangle (1,0.65);
  \fill[blue!16] (2,0) rectangle (3,0.65);
  \fill[blue!28] (3,0) rectangle (4,0.65);
  \node at (0.5,0.33) {强};
  \draw[thick] (3,-0.1) -- (3,0.75); \node[font=\tiny] at (3,-0.30) {小节线};
  \node[anchor=west] at (6.1,0.33) {两小节};
\end{tikzpicture}

\small 图 2-1：4/4 与 3/4 的拍组。颜色深浅只表示强/弱位置，不代表音量必须照做。
\end{center}

\subsection{拍号：每小节几拍、以什么为一拍}

拍号写作分数形式，但不要读成算术分数：

\begin{itemize}
  \item \textbf{分子}：每小节的拍数。
  \item \textbf{分母}：以哪种音符时值为“一拍”，4 表示四分音符，8 表示八分音符，2 表示二分音符。
\end{itemize}

4/4 读作“四四拍”，每小节四拍，以四分音符为一拍；3/4 每小节三拍；2/4 每小节两拍。6/8 是复合拍：每小节六个八分音符，按 3+3 两组组织，重音落在第 1 与第 4 个八分音符上，听起来像“两拍，每拍三下”。

\begin{center}
\small
\begin{tabular}{llll}
\toprule
拍号 & 每小节拍数 & 一拍时值 & 常见重音位置 \\
\midrule
4/4 & 4 & 四分音符 & 1 强，3 次强，2/4 弱 \\
3/4 & 3 & 四分音符 & 1 强，2/3 弱（圆舞曲感） \\
2/4 & 2 & 四分音符 & 1 强，2 弱（进行曲感） \\
6/8 & 6 个八分音符 & 八分音符 & 1 强，4 次强，按 3+3 分组 \\
\bottomrule
\end{tabular}

表 2-1：四个入门拍号。6/8 的“拍”在听觉上常按两组三连音来数。
\end{center}

4/4 与 2/2（cut time，cut common time）都常被叫作“四四拍”的亲戚，但 2/2 是以二分音符为一拍、每小节两拍；听感更快、更宽，写谱时不可与 4/4 混用。

\subsection{音符时值：从全音符到十六分音符}

音符时值用“全音符”作为基准，逐级二分：

\begin{center}
\begin{tikzpicture}[x=1cm,y=1cm,font=\scriptsize]
  \foreach \y/\name/\beats/\sec/\w in {0/全音符/4/2.0/4, 1/二分音符/2/1.0/2, 2/四分音符/1/0.5/1, 3/八分音符/0.5/0.25/0.5, 4/十六分音符/0.25/0.125/0.25} {
    \draw[fill=blue!18,draw=blue!55!black] (0,\y*0.75) rectangle ++(\w,0.38);
    \node[anchor=west] at (4.3,\y*0.75+0.19) {\name：\beats\ 拍，120 BPM 时约 \sec\ 秒};
  }
\end{tikzpicture}

\small 图 2-2：全音符、二分、四分、八分、十六分音符的相对时值（以 120 BPM、四分音符为一拍计）。
\end{center}

附点（dot）把原时值增加一半：附点四分音符 = 1.5 拍，附点二分音符 = 3 拍。连音线（tie）把相邻同音高的音符连成一个更长的音，不重新起音；圆滑线（slur）连接不同音高，表示乐句而不是时值相加——两者外形相近，作用不同。

三连音（triplet）把“两拍的空间”均分成三份，每份 2/3 拍；它是节奏上的“三对二”，不是把拍号改成分数。入门阶段先掌握二分、四分、八分与附点，三连音在第 7 章旋律写作再展开。

\subsection{休止、弱位与弱起}

休止（rest）是“没有声音的时值”，不是“没有音乐”。全休止、二分休止、四分休止、八分休止与同名音符等长；写谱时休止符也算进小节总时值，不能因为“没声音”就漏掉。

弱位（weak position）指每拍内部不被重音强调的位置；弱起（anacrusis，pickup）指段落从弱位开始、第一小节不完整。弱起不是错误：它让段落“先走一步再落地”，流行歌的副歌常用。

判断强弱位不要只看“第几拍”，还要看拍内细分：4/4 的第 1 拍是强位，第 1 拍的后半拍（第 2 个八分）是弱位；切分正是把重音从强位挪到这些弱位上。

\subsection{切分与反拍}

切分（syncopation）是把重音放到弱拍或弱位，让听者感觉“拍子还在，重音不在正拍”。最常见的两种：

\begin{itemize}
  \item \textbf{反拍（backbeat）}：4/4 中把 2、4 拍加重，1、3 拍相对轻；流行鼓的军鼓常打反拍。
  \item \textbf{弱位切分}：在每拍后半拍起音，或让一个音跨过下一拍的正拍，制造“抢拍/拖拍”感。
\end{itemize}

\begin{center}
\begin{tikzpicture}[x=1cm,y=1cm,font=\scriptsize]
  \foreach \i in {0,...,7} {\draw[black!25] (\i,1) rectangle ++(1,0.6);}
  \foreach \i in {0,2,4,6} {\fill[blue!45] (\i+0.1,1.08) rectangle ++(0.8,0.44);}
  \node[anchor=west] at (0,1.95) {\textbf{直拍}};

  \foreach \i in {0,...,7} {\draw[black!25] (\i,0) rectangle ++(1,0.6);}
  \foreach \i in {0,2,3,5,7} {\fill[red!45] (\i+0.1,0.08) rectangle ++(0.8,0.44);}
  \node[anchor=west] at (0,0.80) {\textbf{切分}};
\end{tikzpicture}

\small 图 2-3：同一小节 4/4 的直拍与切分。上排直拍：重音在 1、3、5、7 个八分位置；下排切分：重音移到 3、6、8 等弱位。格子是八个八分音符；色块是起音位置，不代表每个音都一样长。
\end{center}

切分不是“随便错拍”。判断标准是：去掉重音后， underlying pulse（底层的拍）仍然可数；听众能感到“拍在，重音偏了”，才是切分；如果连拍都数不出来，多半是节奏没对齐。

\subsection{速度、律动与风格}

同样的音符，速度快慢会改变律动（groove）。120 BPM 的 4/4 听起来平稳，140 BPM 更急，70 BPM 更拖。风格上的 swing（摇摆）把八分音符奏成“长–短”不等值；straight（直）则保持均分。swing 是演奏/编曲处理，不是拍号变化；写谱时仍记成八分音符，由演奏法或后期处理体现。

入门阶段不要同时改速度、拍号与切分。一次只动一个变量：先把直拍数稳，再加反拍，最后才做弱位切分；这样耳朵能分辨每个变化来自哪里。

\subsection{在夜半里怎么做}

\begin{enumerate}
  \item 新建工程时，在启动页的“详情”里填速度（BPM）、拍号与采样率；拍号分子可步进 1–32，分母可选 1/2/4/8/16/32。工程建好后，标尺按工程拍号画出小节与拍线。
  \item 走带栏有 BPM 显示与输入；改速度后，标尺的小节宽度与播放速度一起变。先在 90–120 BPM 之间选一个能数清的速度。
  \item 菜单里有“节拍器开关”。打开后跟着 click 数拍；4/4 每小节四下，3/4 三下，6/8 按两组三下数。先只用耳朵跟，不看屏幕。
  \item 在编曲区看标尺：竖线是小节线，数字是小节号；放大后能看到拍线与拍内网格。写片段前先确认“这一小节从哪到哪”。
  \item 在钢琴卷帘里用绘制工具（数字键 \texttt{2}）落音符；网格按当前量化值吸附，写快了不会歪。需要修手弹的毛边时，用“编辑 → 音符量化（Quantize 1/16）…”。
  \item 按 \texttt{Space} 播放，跟着走带数拍；如果重音总落在反拍，先检查是不是把音符放在了每拍的后半格，而不是先怀疑监听。
\end{enumerate}

\screenshot{piano-roll-rhythm-grid}{钢琴卷帘节奏网格与拍号}

\textbf{范围说明：}本册把 swing 当作律动与写作概念；夜半鼓生成器里的 swing 参数属于乐器/编曲分册，不在本章展开。拍号、BPM 与网格的操作仍以本章为准。

\subsection{练习}

\begin{enumerate}
  \item 120 BPM、4/4 的一小节多少秒？90 BPM 时呢？\\
  \emph{提示：}一拍 60/BPM 秒，4/4 四拍。\\
  \emph{答案：}120 BPM 为 2 秒；90 BPM 为 \(4 \times 60/90 \approx 2.67\) 秒。
  \item 3/4 与 6/8 都常听起来“三下一组”，它们的小节拍数有什么不同？\\
  \emph{答案：}3/4 每小节三拍、以四分音符为一拍；6/8 每小节六个八分音符，按 3+3 两组组织，重音在第 1、4 个八分。
  \item 附点四分音符等于几拍？附点二分呢？\\
  \emph{答案：}1.5 拍；3 拍。
  \item 把 4/4 的直拍（每拍正拍起音）改成反拍，重音应落在哪些拍？\\
  \emph{答案：}2、4 拍（每小节第二、第四拍）。
  \item 在夜半里建一个 4/4、100 BPM 的工程，打开节拍器，用八分音符写“强–弱–强–弱”两小节；再把第二小节第 2 拍的后半拍加一个音，听切分感。\\
  \emph{提示：}先对准小节线，再数每拍的八个半格；不要用拖动去猜位置。
\end{enumerate}

\subsection{帮助主题映射}

\begin{center}
\footnotesize
\setlength{\tabcolsep}{3pt}
\begin{longtable}{p{0.30\linewidth}p{0.48\linewidth}p{0.16\linewidth}}
\toprule
本章小节 & 应用内帮助锚点 & 备注 \\
\midrule
\endhead
脉冲、拍与速度 & \texttt{music-theory-pulse-beat-tempo} & 概念/式 \eqref{eq:beat-duration} \\
拍号 & \texttt{music-theory-time-signature} & 表 2-1 \\
音符时值 & \texttt{music-theory-note-values} & 图 2-2 \\
休止、弱位与弱起 & \texttt{music-theory-rests-weak-beats} & 概念 \\
切分与反拍 & \texttt{music-theory-syncopation-backbeat} & 图 2-3 \\
速度与律动 & \texttt{music-theory-groove-tempo} & 概念 \\
在夜半里怎么做 & \texttt{music-theory-timeline-grid} & 操作 \\
练习 & \texttt{music-theory-ch2-practice} & 答案在章末 \\
\bottomrule
\end{longtable}
\end{center}
